\section{Open Questions}

The following five questions remain genuinely open after this replication.
Each is scoped so that the ``next steps'' can be executed on free
endpoints (local Stim + numpy on CherryRd / m1) without paid API calls.

\begin{enumerate}
\item \textbf{Resource-cost comparison against modern FT-EC schemes.}
How does the $[\![4,2,2]\!]$ flag-qubit FT-EC circuit compare
quantitatively (logical error rate, qubit overhead, gate count, wall-time
under noise) to the modern flag-qubit distance-3 protocols of
Chao--Reichardt~2018 and Chamberland--Beverland~2018, and to a
distance-3 rotated surface-code patch, at matched physical error rates in
the $p \in [10^{-3}, 3\!\times\!10^{-2}]$ window?

\emph{Basis:} Our replication verified the paper's internal FT scaling
for $L_a$ but did not benchmark this circuit against later-generation
FT-EC schemes; the paper predates the Chao--Reichardt flag formalism
(2018) and the resource-cost story vs.\ surface code is now the
operationally interesting question for near-term FT experiments.

\emph{Next steps:} Implement Chao--Reichardt flag circuits for
$[\![7,1,3]\!]$ Steane and a $d=3$ surface patch in Stim; run matched-$p$
Monte-Carlo scans ($\geq 10^6$ shots each) using the same
depol2+single noise models already in \texttt{work/ft422\_stim.py}; report
LER, qubit budget, and 2-qubit gate count in a common table; free-endpoint
(local Stim) only.

\item \textbf{Robustness of the $p^2$ scaling under hardware-aware noise.}
Does the reproduced $L_a \sim p^2$ scaling exponent
(2.17 Sz / 1.92 Sx) survive under a hardware-aware ion-trap noise model
with leakage, motional-mode-mediated correlated 2-qubit dephasing, and
asymmetric SPAM, or does it collapse toward $p^1$ as the paper's residual
$L_b$ gap suggests would happen?

\emph{Basis:} Our depol2/single models under-predict $L_b$ by
${\sim}2\times$ relative to paper, indicating the ion-trap experiment
carries non-Pauli noise that our sim omits. Whether that noise preserves
FT for $L_a$ (only inflating $L_b$) or degrades $L_a$'s scaling is an open
experimental-vs-theory question that a richer noise model can settle in
simulation.

\emph{Next steps:} Add a Kraus-channel leakage model (three-level qutrit
projection to computational subspace, ${\sim}10^{-4}$ leakage per CNOT) and
correlated $ZZ$-crosstalk term (${\sim}10^{-3}$) as post-CNOT operations
in Stim's TableauSimulator; scan $p$; refit log-log slopes; compare against
the current 2.17/1.92 baseline; free-endpoint (local Stim + numpy).

\item \textbf{Concatenation to $[\![4,2,2]\!] \times [\![4,2,2]\!]$.}
Can the $[\![4,2,2]\!]$ FT-EC building block be extended to a concatenated
$[\![4,2,2]\!]\times[\![4,2,2]\!]$ scheme
(16 physical $\rightarrow$ 4 logical, distance 4) with computable
pseudo-threshold, and how does that pseudo-threshold compare to the
$[\![7,1,3]\!]$ Steane and to surface-code $d=3$?

\emph{Basis:} The paper deliberately targets a single $[\![4,2,2]\!]$
block for near-term hardware; whether concatenation yields a useful FT
primitive at intermediate scale (before full surface-code) is unresolved
and directly relevant to the current 5--50 qubit hardware regime the paper
occupies.

\emph{Next steps:} Construct the concatenated encoding in Stim by nesting
two levels of flag-qubit encoding; do exhaustive weight-2 fault enumeration
to prove $d=4$ detection; run Monte-Carlo LER scan across
$p \in [10^{-4}, 3\!\times\!10^{-2}]$; extract the crossing point vs.\
bare qubit as the pseudo-threshold; free-endpoint (local Stim).

\item \textbf{Gate-order sensitivity of the flag encoding.}
How sensitive is the fault-tolerance verdict of the flag encoding to the
specific order of the four data-CNOTs between the initialization block and
the flag-off gate, i.e., are there permutations of the encoding gate order
that break the 0/324 undetected-$L_a$ property?

\emph{Basis:} Our exhaustive-enumeration test uses one canonical gate
ordering ($q_0\!\to\!\text{flag}$, $q_0\!\to\!q_1$, $q_0\!\to\!q_2$,
$q_0\!\to\!q_3$, $q_0\!\to\!\text{flag}$). The 5-CNOT canonical structure
has $4! = 24$ permutations of the data-CNOT block; the paper does not
enumerate which orderings preserve FT, and a compiler that reorders for
scheduling reasons could silently break FT.

\emph{Next steps:} Wrap the \texttt{ft\_flag\_test.py} enumeration in a
permutation loop over the 24 data-CNOT orderings; count undetected
$L_a$-error faults for each; report which orderings preserve 0 undetected
$L_a$ errors and which do not; free-endpoint (local Stim, ${<}1$ minute
wall).

\item \textbf{Multi-round cycle behavior.}
For a specific near-term ion-trap or superconducting-qubit cycle
experiment (e.g., a 3-round syndrome cycle on the same $[\![4,2,2]\!]$
code with mid-circuit measurement + reset), does the FT property
demonstrated for a single stabilizer measurement round survive across
multiple rounds, and where does correlated-measurement error break the
$p^2$ scaling?

\emph{Basis:} The paper only measures the stabilizer once; real FT-EC
use requires repeated syndrome extraction with mid-circuit reset. Whether
the 0/324 single-fault result extends to 3-round enumeration (which grows
to ${\sim}10^4$ fault points) is not tested here and is the pragmatic
next question for hardware.

\emph{Next steps:} Extend \texttt{work/ft\_flag\_test.py} to a 3-round
syndrome-extraction circuit with mid-circuit measurement + reset in Stim;
do exhaustive single-fault enumeration; report undetected $L_a$-error
count per round; run matched Monte-Carlo LER scan on the multi-round
circuit; free-endpoint (local Stim).
\end{enumerate}
